\section{Síntesis y ampliación}

\subsection{Tabla de síntesis central}

\begin{table}[H]
\centering
\begin{tabular}{p{0.26\textwidth} p{0.24\textwidth} p{0.25\textwidth}}
\toprule
Tema & Conclusión central & Consejo de ingeniería \\
\midrule
Núcleo de Q-Learning & Se aprende la función Q óptima; la política se obtiene mediante $\arg\max_a Q(s,a)$ & Prepara el replay buffer, la target network y la regularization \\
DQN & replay buffer + target network + $\epsilon$-greedy es la versión escalable de Q-Learning & Mantén el buffer size, la diversidad de muestreo y la frecuencia de actualización de la target network \\
Estabilización & Double Q, la prevención del reward hack y el reward shaping & Monitorea el TD error y la reward variance; configura un rollback guardrail \\
Despliegue de ingeniería & Matriz de evaluación + checklist + telemetry & Verifica el checkpoint con sim rollout + human-in-the-loop replay; monitorea el Q drift con un guardrail en producción \\
\bottomrule
\end{tabular}
\caption{Correspondencia entre método e ingeniería en Lecture 6}
\end{table}

\subsection{Lecturas adicionales}
\subsection{Lecturas adicionales}

\begin{itemize}
    \item Mnih et al., \textit{Playing Atari with Deep Reinforcement Learning} (DQN, 2013)
    \item Van Hasselt et al., \textit{Deep Reinforcement Learning with Double Q-learning} (Double DQN, 2016)
    \item Lillicrap et al., \textit{Continuous Control with Deep Reinforcement Learning} (DDPG, 2016)
    \item Schaul et al., \textit{Prioritized Experience Replay} (ICLR 2016)
    \item Jaderberg et al., \textit{Population Based Training of Neural Networks} (NeurIPS 2017) — for lifecycle automation
\end{itemize}

\subsection{Resumen de la sección}

El valor de Q-Learning radica en derivar la política a partir de la función de valor; al combinarlo con la estabilización, la gobernanza de datos y el despliegue de ingeniería, este algoritmo clásico puede operar de manera confiable en los sistemas de RL modernos.

\subsection{Recordatorios de acción}

\begin{itemize}
    \item Después de cada ronda de entrenamiento, documenta la evidence matrix y el order-of-operations para facilitar el próximo rollback.
    \item Monitorea `TD error drift`, `reward variance` y `episode success rate`; en cuanto cualquiera de estos indicadores active el guardrail, toma un snapshot de inmediato.
    \item Programa una quarterly replay audit (human review de 20 failure trajectories).
\end{itemize}

\end{document}
